\section{Dense thresholds, matched certification and precision requirements}\label{sec:s10}
\subsection{Experiments fixed before the additional runs}
The additional protocol contains three modules: six pairing-ablation cells with 500 replications each, six dense-grid cells with 200 replications each, and seven robustness cells with 300 replications each. This gives 6,300 design--grid replications and 37,800 method records. The dense-grid cells deliberately reuse each generated dataset at three grid resolutions: they contain 400 distinct datasets, not 1,200 independent datasets. The additional modules together use 5,500 distinct generated datasets. All six methods within a replication share training, reference masks, gate and evaluation data. These runs are separate from the 6,000 primary replications.

The ablation repeats alignment at reference fractions $0.05,0.10,0.20$ and weak, reversed and absent signals at fraction $0.10$, always with $n=6,000$, $h=0.35$ and seven thresholds. Unpaired uses the absolute-critical-value bounds of Section S3.5, with statistical errors $0.0245$ for the anchor and for the candidate maximum. The two upper simulated quantiles each receive error $0.00025$. Its Gaussian numerical intervals are exactly the same intervals used by Paired and have total error $0.0005$. Thus each selector has total statistical-plus-numerical error $0.05$. Each absolute Taylor expansion receives the same $b_N$ cushion; the sum is $2b_N$, compared with Paired's $\lambda b_N$. This is part of the error-cancellation comparison, not a different reference budget. The other four methods are Anchor, Full, Trace and Plug-in. The matrix Projection comparator remains in the primary experiment.

The dense experiments use $J=21,51,101$ thresholds $0.8\Phi^{-1}(u_j)$ with evenly spaced $u_j\in[0.02,0.98]$. The 21- and 51-point grids are subsets of the 101-point grid. Both alignment with random verification and local reversal with selective verification use $n=6,000$, $p=0.10$ and $h=0.35$. A common enlargement $10^{-8}I_J$ is declared for these calculations. Every reference mask, candidate regression and generated dataset is shared across grid resolutions. We evaluate the raw simultaneous band on its grid and its monotone envelope on a common 1,001-point grid. The latter is a numerical check of the envelope; the all-threshold guarantee follows from monotonicity as proved in Section S4. The experiment does not replace a growing-grid regularity condition by an empirical coverage estimate.

\subsection{Outcome and nuisance-regression robustness}
The three continuous non-Gaussian designs replace the normal error by a standardised $t_3$, $(\chi^2_3-3)/\sqrt6$, or $(A+0.5Z)/\sqrt{1.25}$ with independent $A\in\{-1,1\}$ equally likely and $Z\sim N(0,1)$. Mean and scale functions are unchanged. These errors have mean zero and variance one. Candidate Gaussian distribution regressions can be misspecified, while their threshold predictions remain bounded. Surrogate noise is Gaussian with standard deviation $0.4$.

For the ordinal design, generate the normal latent outcome and discretise it using nine cutpoints equally spaced from $-2$ to $2$, producing categories $0,\ldots,9$. Set $S=Y+0.4\eta$ and evaluate all nine nontrivial cumulative probabilities. True probabilities use the latent normal distribution at the category cutpoints. The two-dimensional design takes independent $X_1,X_2\sim U[-1,1]$, mean $1.2X_1+0.6X_1^2+0.8X_2$, scale $0.8+0.25(X_1^2+X_2^2)$, a product Epanechnikov kernel with $h=0.5$, and target $(0,0)$. Its local information is $n_Gph^2$.

A separate nonparametric-candidate design uses empirical neighbour distribution functions: the anchor searches in $X$, and the surrogate candidate in $(X,S)$. Coordinates are divided by their verified-training standard deviations, floored at $0.1$. The neighbour count is the smaller of the verified-training count and $\max\{20,\lfloor\sqrt{n_{T,R}}\rfloor\}$. All nine nontrivial category thresholds are used for the ordinal design; other robustness cells have 21 thresholds on the same probability range as the dense-grid study. All have $n=6,000$, $p=0.10$, and common enlargement $10^{-8}I_J$.

The learned-verification sensitivity uses
$\pi(X)=\operatorname{expit}(a+0.8X)$, with $a$ determined by $\E\pi(X)=0.10$. A logistic regression of $R$ on $(1,X)$ is fitted only in the training role; its predictions are clipped to $[0.005,0.5]$. These fitted probabilities replace the known design probabilities in the downstream signals. This cell measures actual plug-in sensitivity and is reported separately from the known-probability calibration result. Learned-probability inference requires the local drift and score-replacement conditions in Section S8; fitting a correct propensity model alone does not establish those product-rate conditions for every misspecified outcome regression.

For each design, the reference CDF is computed by integrating its known conditional distribution over the kernel neighbourhood, using 100-node Gauss--Legendre quadrature per covariate. Independent refinement to 180 nodes is checked. Population critical-value gains and harmful-selection probabilities are reported only in the normal, one-dimensional, Gaussian-candidate known-design cases where the conditional Gaussian covariance calculation applies. They are left unavailable in the other cells, rather than replacing an exact population comparison by an unlabelled finite truth sample. Coverage, width, distribution loss and adoption remain observed in every cell.

\subsection{Numerical refinement at dense thresholds}
The first five replication identifiers in each dense-grid cell are re-evaluated with 8,192 and 32,768 radial directions. These 30 comparisons were fixed before the refinement. Both calculations retain the same numerical critical-value intervals and total confidence budgets. We report changes in critical values, score radii, chosen weights, evaluation widths and point estimates. Weight agreement alone is not a measure of Gaussian integration accuracy when several grid choices have similar lower gains. The delivered records retain every disagreement and its inferential consequences. Independent directional finite differences also check the gradient in dimensions 21, 51 and 101. These calculations validate the implementation of the numerical derivative and measure integration sensitivity; the separate direct-Gaussian order-statistic construction supplies the numerical critical-value probability guarantee.

\subsection{Additional results and interpretation}
The completed run contains 6,300 design--grid replications, with no execution failures. Table~\ref{tab:s10-validation} gives coverage, width and adoption for every additional cell and method. Exact binomial intervals and Monte Carlo standard errors for all continuous summaries accompany the machine-readable table. These are cellwise Monte Carlo intervals, not a simultaneous assertion across all simulation cells. The dense output eigenvalues and local reference counts are retained alongside each replication. The 1,001-point envelope check is separate from grid coverage.

In the paired ablation, width differences and their standard errors are computed within the same generated sample. Two absolute critical-value bounds can remain too wide to certify a gain even when the paired difference can do so. The lower-bound constants and the allocation of error are fully disclosed in Section S3.5; this experiment isolates the stated implementations rather than every possible unpaired optimisation.

In the 30 dense refinement cases, 25 weights agree. Maximum absolute changes in radial critical values and paired radial gains are 0.005813 and 0.004265. Maximum relative changes in an evaluated band width and in the corresponding population critical value are 3.258\% and 2.140\%. No refined selected candidate in these cases has a negative assessed gain. Selection labels, integration error and inferential impact are therefore all recorded, rather than requiring equality of a mesh index as a proxy for numerical accuracy.

The seven robustness cells retain Gaussian-candidate misspecification or use the declared neighbour estimator. Their true CDFs are evaluated analytically in the outcome dimension and by covariate quadrature; the independent 180-node refinement and implementation tests are archived. These checks do not convert an asymptotic statement into a finite-sample guarantee for an estimated verification model.

\subsection{Extended EPA precision experiment}
The original 1,200 budget--design--mask combinations and 126,000 threshold records remain unchanged. The added random-audit grid has 1,600 budget--mask combinations, 24,000 method--humidity records and 168,000 thresholdwise intervals. It uses the unchanged historical candidates and the same mask generator as the original study. The first 200 uniform-draw families are nested with those already reported; the next 200 are additional families. No thresholds are selected by their sign-resolution rates. Table~\ref{tab:s10-epa} reports all methods on the added budgets. Complete targetwise coverage, false-sign family probabilities, adoption and interval-width summaries are in \texttt{results/v3/reporting}.

The first-tested-budget summary distinguishes an observed resolution rate above 0.80 from an exact 95\% Monte Carlo lower bound above 0.80. It does not interpolate a minimum required budget. For $p=1$ the entire audit frame is observed and the contrast is known; this exact census endpoint is not passed through a Gaussian covariance enlargement or represented as an additional simulation. Historical calibration measurements remain an additional cost at every budget.
\begin{longtable}{llrrrr}
\caption{Complete additional validation. Coverage and adoption are frequencies; width is the mean within-replication ratio to Anchor.}\label{tab:s10-validation}\\
\toprule Cell & Method & Coverage & MCSE & Width & Adoption \\\midrule\endfirsthead
\toprule Cell & Method & Coverage & MCSE & Width & Adoption \\\midrule\endhead
\bottomrule\endfoot
\multicolumn{6}{l}{101: aligned, $p=0.05$}\\
101 & Anchor & 0.954 & 0.009 & 1.000 & 0.000 \\
101 & Full & 0.966 & 0.008 & 0.723 & 1.000 \\
101 & Trace & 0.968 & 0.008 & 0.722 & 1.000 \\
101 & Plug-in & 0.968 & 0.008 & 0.724 & 1.000 \\
101 & Paired & 0.952 & 0.010 & 0.802 & 0.858 \\
101 & Unpaired & 0.954 & 0.009 & 1.000 & 0.000 \\
\multicolumn{6}{l}{102: aligned, $p=0.10$}\\
102 & Anchor & 0.948 & 0.010 & 1.000 & 0.000 \\
102 & Full & 0.952 & 0.010 & 0.724 & 1.000 \\
102 & Trace & 0.954 & 0.009 & 0.724 & 1.000 \\
102 & Plug-in & 0.958 & 0.009 & 0.724 & 1.000 \\
102 & Paired & 0.954 & 0.009 & 0.747 & 0.998 \\
102 & Unpaired & 0.948 & 0.010 & 1.000 & 0.000 \\
\multicolumn{6}{l}{103: aligned, $p=0.20$}\\
103 & Anchor & 0.958 & 0.009 & 1.000 & 0.000 \\
103 & Full & 0.954 & 0.009 & 0.759 & 1.000 \\
103 & Trace & 0.956 & 0.009 & 0.759 & 1.000 \\
103 & Plug-in & 0.954 & 0.009 & 0.759 & 1.000 \\
103 & Paired & 0.956 & 0.009 & 0.769 & 1.000 \\
103 & Unpaired & 0.958 & 0.009 & 0.993 & 0.032 \\
\multicolumn{6}{l}{104: weak, $p=0.10$}\\
104 & Anchor & 0.928 & 0.012 & 1.000 & 0.000 \\
104 & Full & 0.938 & 0.011 & 0.956 & 1.000 \\
104 & Trace & 0.938 & 0.011 & 0.957 & 0.998 \\
104 & Plug-in & 0.934 & 0.011 & 0.957 & 0.998 \\
104 & Paired & 0.928 & 0.012 & 1.000 & 0.002 \\
104 & Unpaired & 0.928 & 0.012 & 1.000 & 0.000 \\
\multicolumn{6}{l}{105: reversal, $p=0.10$}\\
105 & Anchor & 0.942 & 0.010 & 1.000 & 0.000 \\
105 & Full & 0.946 & 0.010 & 1.110 & 1.000 \\
105 & Trace & 0.942 & 0.010 & 1.000 & 0.004 \\
105 & Plug-in & 0.942 & 0.010 & 1.000 & 0.004 \\
105 & Paired & 0.942 & 0.010 & 1.000 & 0.000 \\
105 & Unpaired & 0.942 & 0.010 & 1.000 & 0.000 \\
\multicolumn{6}{l}{106: no signal, $p=0.10$}\\
106 & Anchor & 0.958 & 0.009 & 1.000 & 0.000 \\
106 & Full & 0.958 & 0.009 & 1.002 & 1.000 \\
106 & Trace & 0.958 & 0.009 & 1.000 & 0.496 \\
106 & Plug-in & 0.958 & 0.009 & 1.000 & 0.484 \\
106 & Paired & 0.958 & 0.009 & 1.000 & 0.000 \\
106 & Unpaired & 0.958 & 0.009 & 1.000 & 0.000 \\
\multicolumn{6}{l}{206: aligned, $J=21$}\\
206 & Anchor & 0.960 & 0.014 & 1.000 & 0.000 \\
206 & Full & 0.940 & 0.017 & 0.740 & 1.000 \\
206 & Trace & 0.940 & 0.017 & 0.740 & 1.000 \\
206 & Plug-in & 0.940 & 0.017 & 0.741 & 1.000 \\
206 & Paired & 0.955 & 0.015 & 0.765 & 1.000 \\
206 & Unpaired & 0.960 & 0.014 & 1.000 & 0.000 \\
\multicolumn{6}{l}{207: aligned, $J=51$}\\
207 & Anchor & 0.960 & 0.014 & 1.000 & 0.000 \\
207 & Full & 0.965 & 0.013 & 0.750 & 1.000 \\
207 & Trace & 0.960 & 0.014 & 0.750 & 1.000 \\
207 & Plug-in & 0.965 & 0.013 & 0.751 & 1.000 \\
207 & Paired & 0.965 & 0.013 & 0.782 & 0.985 \\
207 & Unpaired & 0.960 & 0.014 & 1.000 & 0.000 \\
\multicolumn{6}{l}{208: aligned, $J=101$}\\
208 & Anchor & 0.960 & 0.014 & 1.000 & 0.000 \\
208 & Full & 0.960 & 0.014 & 0.753 & 1.000 \\
208 & Trace & 0.960 & 0.014 & 0.753 & 1.000 \\
208 & Plug-in & 0.960 & 0.014 & 0.754 & 1.000 \\
208 & Paired & 0.955 & 0.015 & 0.798 & 0.965 \\
208 & Unpaired & 0.960 & 0.014 & 1.000 & 0.000 \\
\multicolumn{6}{l}{209: reversal, $J=21$}\\
209 & Anchor & 0.935 & 0.017 & 1.000 & 0.000 \\
209 & Full & 0.925 & 0.019 & 1.128 & 1.000 \\
209 & Trace & 0.935 & 0.017 & 1.000 & 0.010 \\
209 & Plug-in & 0.935 & 0.017 & 1.000 & 0.005 \\
209 & Paired & 0.935 & 0.017 & 1.000 & 0.000 \\
209 & Unpaired & 0.935 & 0.017 & 1.000 & 0.000 \\
\multicolumn{6}{l}{210: reversal, $J=51$}\\
210 & Anchor & 0.945 & 0.016 & 1.000 & 0.000 \\
210 & Full & 0.940 & 0.017 & 1.124 & 1.000 \\
210 & Trace & 0.945 & 0.016 & 1.000 & 0.010 \\
210 & Plug-in & 0.945 & 0.016 & 1.000 & 0.010 \\
210 & Paired & 0.945 & 0.016 & 1.000 & 0.000 \\
210 & Unpaired & 0.945 & 0.016 & 1.000 & 0.000 \\
\multicolumn{6}{l}{211: reversal, $J=101$}\\
211 & Anchor & 0.950 & 0.015 & 1.000 & 0.000 \\
211 & Full & 0.930 & 0.018 & 1.123 & 1.000 \\
211 & Trace & 0.950 & 0.015 & 1.000 & 0.010 \\
211 & Plug-in & 0.950 & 0.015 & 1.000 & 0.010 \\
211 & Paired & 0.950 & 0.015 & 1.000 & 0.000 \\
211 & Unpaired & 0.950 & 0.015 & 1.000 & 0.000 \\
\multicolumn{6}{l}{301: t3, $d=1$, gaussian, random}\\
301 & Anchor & 0.967 & 0.010 & 1.000 & 0.000 \\
301 & Full & 0.960 & 0.011 & 0.809 & 1.000 \\
301 & Trace & 0.953 & 0.012 & 0.807 & 1.000 \\
301 & Plug-in & 0.960 & 0.011 & 0.807 & 1.000 \\
301 & Paired & 0.963 & 0.011 & 0.841 & 0.943 \\
301 & Unpaired & 0.967 & 0.010 & 1.000 & 0.000 \\
\multicolumn{6}{l}{302: skew, $d=1$, gaussian, random}\\
302 & Anchor & 0.960 & 0.011 & 1.000 & 0.000 \\
302 & Full & 0.960 & 0.011 & 0.779 & 1.000 \\
302 & Trace & 0.960 & 0.011 & 0.778 & 1.000 \\
302 & Plug-in & 0.960 & 0.011 & 0.779 & 1.000 \\
302 & Paired & 0.960 & 0.011 & 0.807 & 0.980 \\
302 & Unpaired & 0.960 & 0.011 & 1.000 & 0.000 \\
\multicolumn{6}{l}{303: mixture, $d=1$, gaussian, random}\\
303 & Anchor & 0.953 & 0.012 & 1.000 & 0.000 \\
303 & Full & 0.977 & 0.009 & 0.723 & 1.000 \\
303 & Trace & 0.977 & 0.009 & 0.723 & 1.000 \\
303 & Plug-in & 0.977 & 0.009 & 0.723 & 1.000 \\
303 & Paired & 0.970 & 0.010 & 0.744 & 1.000 \\
303 & Unpaired & 0.953 & 0.012 & 1.000 & 0.000 \\
\multicolumn{6}{l}{304: discrete, $d=1$, gaussian, random}\\
304 & Anchor & 0.947 & 0.013 & 1.000 & 0.000 \\
304 & Full & 0.963 & 0.011 & 0.667 & 1.000 \\
304 & Trace & 0.960 & 0.011 & 0.665 & 1.000 \\
304 & Plug-in & 0.960 & 0.011 & 0.666 & 1.000 \\
304 & Paired & 0.953 & 0.012 & 0.686 & 1.000 \\
304 & Unpaired & 0.950 & 0.013 & 0.994 & 0.020 \\
\multicolumn{6}{l}{305: normal, $d=2$, gaussian, random}\\
305 & Anchor & 0.937 & 0.014 & 1.000 & 0.000 \\
305 & Full & 0.957 & 0.012 & 0.754 & 1.000 \\
305 & Trace & 0.963 & 0.011 & 0.754 & 1.000 \\
305 & Plug-in & 0.967 & 0.010 & 0.755 & 1.000 \\
305 & Paired & 0.950 & 0.013 & 0.822 & 0.877 \\
305 & Unpaired & 0.937 & 0.014 & 1.000 & 0.000 \\
\multicolumn{6}{l}{306: normal, $d=1$, knn, random}\\
306 & Anchor & 0.983 & 0.007 & 1.000 & 0.000 \\
306 & Full & 0.960 & 0.011 & 0.753 & 1.000 \\
306 & Trace & 0.960 & 0.011 & 0.753 & 1.000 \\
306 & Plug-in & 0.960 & 0.011 & 0.753 & 1.000 \\
306 & Paired & 0.963 & 0.011 & 0.777 & 0.993 \\
306 & Unpaired & 0.983 & 0.007 & 1.000 & 0.000 \\
\multicolumn{6}{l}{307: normal, $d=1$, gaussian, estimated}\\
307 & Anchor & 0.960 & 0.011 & 1.000 & 0.000 \\
307 & Full & 0.970 & 0.010 & 0.737 & 1.000 \\
307 & Trace & 0.970 & 0.010 & 0.737 & 1.000 \\
307 & Plug-in & 0.970 & 0.010 & 0.738 & 1.000 \\
307 & Paired & 0.963 & 0.011 & 0.768 & 0.997 \\
307 & Unpaired & 0.960 & 0.011 & 1.000 & 0.000 \\
\end{longtable}
\begin{longtable}{rlrrr}
\caption{Added random EPA budgets. Coverage is joint over 21 threshold--humidity combinations; width is averaged across humidity within each mask and divided by Anchor. Each row uses 400 masks.}\label{tab:s10-epa}\\
\toprule $p$ (\%) & Method & Joint coverage & MCSE & Width ratio \\\midrule\endfirsthead
\toprule $p$ (\%) & Method & Joint coverage & MCSE & Width ratio \\\midrule\endhead
\bottomrule\endfoot
40 & Anchor & 0.968 & 0.009 & 1.000 \\
40 & Full & 0.973 & 0.008 & 0.654 \\
40 & Trace & 0.975 & 0.008 & 0.654 \\
40 & Plug-in & 0.970 & 0.009 & 0.655 \\
40 & Paired & 0.975 & 0.008 & 0.701 \\
60 & Anchor & 0.963 & 0.009 & 1.000 \\
60 & Full & 0.978 & 0.007 & 0.651 \\
60 & Trace & 0.978 & 0.007 & 0.651 \\
60 & Plug-in & 0.978 & 0.007 & 0.652 \\
60 & Paired & 0.968 & 0.009 & 0.671 \\
80 & Anchor & 0.960 & 0.010 & 1.000 \\
80 & Full & 0.960 & 0.010 & 0.650 \\
80 & Trace & 0.960 & 0.010 & 0.650 \\
80 & Plug-in & 0.960 & 0.010 & 0.651 \\
80 & Paired & 0.960 & 0.010 & 0.662 \\
95 & Anchor & 0.968 & 0.009 & 1.000 \\
95 & Full & 0.945 & 0.011 & 0.650 \\
95 & Trace & 0.943 & 0.012 & 0.650 \\
95 & Plug-in & 0.948 & 0.011 & 0.650 \\
95 & Paired & 0.948 & 0.011 & 0.658 \\
\end{longtable}

